\subsection{Hahn-Banach 定理（几何形式）}

定理 1（Hahn-Banach）. --- 设 A 是拓扑向量空间 E 的一个非空开凸集，M 是一个不与 A 相交的非空线性流形。则存在一个闭超平面 H，它包含 M 且不与 A 相交。

经平移，问题可化为 \(0 \in A\) 的情形，从而 A 是吸收的。设 p 是 A 的度规 (II, 第 20 页)，于是 A 是由那些点 \(x \in E\) （它们使 \(p(x) < 1\) ）所成的集。另一方面，设 V 是由 M 生成的 E 的向量子空间；于是 M 是 V 中一个不含 0 的超平面，因而存在 V 上唯一的线性型 f，使 M 是由那些点 \(y \in V\) （它们满足 \(f(y) = 1\) ）所成的集。于是假设 \(M \cap A = \varnothing\) 蕴含对所有 \(y \in V\) （使 \(f(y) = 1\) ），有 \(p(y) \geqslant 1\) ；由于 f 和 p 都是正齐次的，\(f(y) \leqslant p(y)\) （对所有 \(y \in V\) 满足 \(f(y) > 0\) ） ；最后，由于 \(p(y) \geqslant 0\) （对所有 \(y \in V\) ），我们看出 \(f(y) \leqslant p(y)\) （对所有 \(y \in V\) ）。由 Hahn-Banach 定理的解析形式 (II, 第 22 页, th. 1)，存在 E 上的一个线性型 h，它扩张 f 且使对所有 \(x \in E\) 有 \(h(x) \leqslant p(x)\) 。设 H 是 E 中方程为 \(h(x) = 1\) 的超平面。显然 \(H \cap V = M\) 且 \(H \cap A = \varnothing\) 。另一方面，H 在 E 中的余集包含非空开集 A，因此 H 在 E 中是闭的 (I, 第 11 页, 推论)。

证毕。

注记. --- 1) 当 \(0 \in \mathbf{M}\) 时，th. 1 可陈述如下：存在 E 中的一个连续线性型，使在 M 中 \(g(x) = 0\) 而在 A 中 \(g(x) > 0\) (II, 第 8 页, prop. 4)。

\begin{enumerate}
\def\labelenumi{\arabic{enumi})}
\setcounter{enumi}{1}
\tightlist
\item
  若把定理 1 应用于 E 带有最细局部凸拓扑的情形 (II, 第 25 页, 例 2)，并且为简单起见假设 \(0 \in A\) ，则得到下面的结果（表面上不涉及拓扑）：若 A 是实向量空间 E 中的一个吸收凸集，且 M 是一个不与 A 相交的非空线性流形，则存在一个包含 M 的超平面 H，使 A 位于 H 的一侧。这个结果并非对每个凸集 A 都成立 (II, 第 65 页, exerc. 5)。
\end{enumerate}

